\documentclass[10pt,a4paper]{amsart}
\usepackage[margin=19mm]{geometry}
\usepackage{amsmath,amssymb,mathtools}
\usepackage[scr=rsfs,cal=euler]{mathalfa}
\usepackage{xfrac}
\usepackage[unicode,hidelinks,bookmarks=false]{hyperref}
\newcommand{\ie}{i.e.\ }
\newcommand{\cf}{cf.\ }
\newcommand{\resp}{respectively\ }
\newcommand{\Subsection}{\subsection}
\providecommand{\nobreakdash}{\nobreak\hbox{-}}
\setlength{\emergencystretch}{2em}
\multlinegap0pt
\usepackage{bm}
\usepackage{bbm}
\usepackage{dsfont}
\usepackage{xspace}
\usepackage{cancel}
\usepackage{mathtools}
\usepackage{amsthm}
\makeatletter
\@ifundefined{theorem}{\newtheorem{theorem}{Theorem}}{}
\makeatother
\newcommand{\lk}{\mathop{\mathrm{lk}}\nolimits}
\newcommand{\sign}{\mathop{\mathrm{sign}}\nolimits}
\title[Cosmetic surgeries on knots in homology spheres and the Cass]{Cosmetic surgeries on knots in homology spheres and the Casson--Walker invariant}
\author{Kazuhiro Ichihara, In Dae Jong}
\date{Source: arXiv:2506.02682v2 (CC BY 4.0)}
\begin{document}
\maketitle

\noindent\emph{Source and licence.} Cosmetic surgeries on knots in homology spheres and the Casson--Walker invariant. Version arXiv:2506.02682v2 (CC BY 4.0), distributed under the Creative Commons Attribution 4.0 International licence, as recorded in the arXiv licence metadata for this exact version and shown on the abstract page https://arxiv.org/abs/2506.02682v2. Full rights notice: licenses/arxiv-2506.02682-CC-BY-4.0.txt.\par\smallskip

\noindent\textit{Editorial scope.} Counted unit: the Theorem whose source label is \texttt{None} in the article (source0.tex lines 338--341, source label \texttt{None}), delivered together with the article's own complete original proof of it (source0.tex lines 343--417, one \texttt{proof} environment of 75 lines). No mathematics is written, repaired or re-derived: statement and proof are the article's own text line for line. Required context retained uncounted: eq:Ito (lines 217--223); eq:DedekindSymbol (lines 242--245); CGformula (lines 252--254). Every definition, display and label of those lines is retained unchanged. Omitted from this package and not redistributed: 1--216, 224--241, 246--251, 255--337, 342--342, 418--611. Nothing inside a retained excerpt is omitted or altered: every retained body is delivered exactly as the article's lines are written, so no omission receipt of any kind accompanies this package. Citations: the retained excerpts carry no citation command at all, so no bibliography entry of the article is transcribed and no reference is redistributed. Reference adaptation: none was needed and none was made; every reference inside the delivered unit resolves inside it, so no unresolved reference or citation is produced. Printed slips kept literal: no printed word, symbol, hypothesis or number of the counted statement or of its proof was altered. Layout and class: the article's own class is \texttt{amsart} with packages including graphicx, overpic, xcolor; the wrapper is the lane's pinned \texttt{amsart} editorial wrapper at 10pt on A4, which restates the theorem-like environments the retained text needs and adds the article's own macro definitions that the retained text needs (\texttt{\textbackslash lk}, \texttt{\textbackslash sign}), with extra wrapper packages (bm, bbm, dsfont, xspace, cancel, mathtools). The delivered numbering is the unit's own. Assets: no retained span contains a figure environment, a graphics-inclusion command, a PDF-page inclusion, a table or a TikZ picture; no image file is copied into this package, the package declares no asset dependency and no image file is redistributed with it. Licence: version arXiv:2506.02682v2 of this article is distributed under the Creative Commons Attribution 4.0 International licence, as recorded in the arXiv licence metadata for this exact version and shown on the abstract page https://arxiv.org/abs/2506.02682v2. A full rights notice is delivered at licenses/arxiv-2506.02682-CC-BY-4.0.txt. Characters: every retained excerpt is pure ASCII, so the delivered engine, XeLaTeX with its default Latin Modern text font, reports no missing character.
\begin{align}\label{eq:Ito}
\begin{split}
D \left(\frac{\lambda_w}{2}-\frac{\varsigma}{8}\right) &= a_2(K_x)\frac{p_y}{q_y}-\frac{p_y}{24q_y} -\frac{p_y}{24q_yq_x^2} +\frac{p_y\ell^2}{24q_y}\\
&\quad +a_2(K_y)\frac{p_x}{q_x} -\frac{p_x}{24q_x}-\frac{p_x}{24q_xq_y^2}+\frac{p_x\ell^2}{24q_x}\\
&\quad + 2v_3(L) + \frac{D}{24}\left( S\left(\frac{p_x}{q_x}\right)-\frac{p_x}{q_x} +S\left(\frac{p_y}{q_y}\right)-\frac{p_y}{q_y} \right).
\end{split}
\end{align}

\begin{equation}\label{eq:DedekindSymbol}
    S\left(\frac{p}{q}\right) + S\left(\frac{q}{p}\right) = 
\frac{p}{q} + \frac{q}{p} + \frac{1}{pq} - 3 \sign (p q).
\end{equation}

\begin{equation}\label{CGformula}
\tau(\Sigma_{K}(p/ q)) = -4p \cdot s(q,p) - \sigma(K,p).
\end{equation}

\begin{theorem}
    Let $L = K \cup K_0$ be a 2-component link in $S^3$ with $\lk(K,K_0)=0$. 
    If $a_2 (K) - q_0 a_3(L) \ne 0$, then $L(p/q,1/q_0) \not\cong L(p/q', 1/q_0)$ for a non-zero integer $q_0$ and distinct non-zero integers $q, q'$ coprime to $p>0$. 
\end{theorem}

\begin{proof}
Suppose that $L(p/q,1/q_0) \cong L(p/q', 1/q_0)$, and set $\lambda_w = \lambda_w ( L(p/q,1/q_0))$, $\lambda_w' = \lambda_w ( L(p/q', 1/q_0))$. 
Then, $\lambda_w = \lambda_w'$ holds, and we will deduce that $a_2 (K) - q_0 a_3(L) = 0$ holds. 
To do so, we apply the formula~\eqref{eq:Ito} in two ways; 
apply~\eqref{eq:Ito} with $p_x=p$, $q_x=q$, $p_y=1$, $q_y=q_0$, and $p_x=p$, $q_x=q'$, $p_y=1$, $q_y=q_0$. 

By the assumption $\ell=0$, the linking matrices, $A$ and $A'$, of the two manifolds $L(p/1,1/q_0)$ and $L(p/q', 1/q_0)$ are given by
\[ A = \begin{pmatrix} \frac{p}{q} & 0 \\ 0 & \frac{1}{q_0} \end{pmatrix} \text{ and } A'=\begin{pmatrix} \frac{p}{q'} & 0 \\ 0 & \frac{1}{q_0} \end{pmatrix}, \] 
respectively. 
Thus, their determinants, $D$ and $D'$, are given by 
\[ D = \frac{p}{q\,q_0} \text{ and } D'=\frac{p}{q'\,q_0} , \]
respectively. 
Also, for their signatures, we have
\[
\varsigma = 
\begin{cases}
    -2 & \text{if } q,q_0 <0 \\
    0 & \text{if } q q_0 <0 \\
    2 & \text{if } q,q_0 >0
\end{cases}
\]
and the same holds for $\varsigma'$. 
Thus, it follows that 
\[ \varsigma - \varsigma' = \mathrm{sign}(q) - \mathrm{sign}(q').\]

Then, applying the formula~\eqref{eq:Ito} in two ways, and taking their difference, we obtain the following. 
\begin{align}\label{Eq2}
\begin{split}
0
&= \frac{\varsigma - \varsigma'}{8} 
+ \frac{q-q'}{p} \left( a_2(K) + 2 q_0 \,v_3(L) \right)
+ \frac{q-q'}{24p} \left( \frac{1}{qq'} - 1 \right)\\
&\quad + \frac{1}{24} \left( \left( S\left( \frac{p}{q} \right) - \frac{p}{q} \right) - \left( S\left( \frac{p}{q'} \right) - \frac{p}{q'} \right) \right) . 
\end{split}
\end{align}

By $1/q_0$-surgery on $K_0$ in $S^3$, we have an integral homology 3-sphere $\Sigma$. 
Then, the knot $K$ can be regarded as a knot in $\Sigma$ naturally. 
Let $K(p/q), K(p/q')$ denote the 3-manifold obtained from $\Sigma$ by $p/q$- and $p/q'$-surgeries on $K$, respectively. 
That is, $L(p/q,1/q_0) \cong K(p/q)$ and $L(p/q', 1/q_0) \cong K(p/q')$ hold. 
This implies that the assumption $L(p/q,1/q_0) \cong L(p/q', 1/q_0)$ is equivalent to $K(p/q) \cong K(p/q')$. 
Since $\Sigma$ is an integral homology sphere, using the surgery formula of the Casson--Gordon invariant (Equation~\eqref{CGformula}), together with $p>0$, we have 
\[
\tau (K(p/q)) - \tau (K(p/q')) = 4p\left( s(q',p) - s(q,p) \right)= 0. 
\]
On the other hand, from Equation~\eqref{eq:DedekindSymbol}, we have 
\begin{align*}
S\left(\frac{p}{q}\right) 
&= \frac{p}{q} + \frac{q}{p} + \frac{1}{pq} - 3 \mathrm{sign}(pq) - S \left( \frac{q}{p} \right) \\
&= 
\frac{p}{q} + \frac{q}{p} + \frac{1}{pq} - 3 \mathrm{sign}(pq) - 
12 \mathrm{sign}(p) s(q,p) 
\end{align*} 
and 
\[
S\left(\frac{p}{q'}\right) 
= 
\frac{p}{q'} + \frac{q'}{p} + \frac{1}{pq'} - 3 \mathrm{sign}(pq') - 
12 \mathrm{sign}(p) s(q',p).
\]
Therefore, 
it follows that 
\[
\left( S\left( \frac{p}{q} \right) - \frac{p}{q} \right)
- \left( S\left( \frac{p}{q'} \right) - \frac{p}{q'} \right) = 
\frac{q-q'}{p}\left( 1 - \frac{1}{qq'} \right) - 3 \left( \mathrm{sign}(q) - \mathrm{sign}(q') \right).
\]
Then, it follows from Equation~\eqref{Eq2} that 
\[
0
 = \frac{ q - q' }{ p } \left( a_2 (K) + 2 q_0 v_3(L) \right) 
 = \frac{ q - q' }{ p } \left( a_2 (K) - q_0 a_3(L) \right).
 \]
Consequently, it is shown that $a_2 (K) = q_0 a_3(L)$ holds. 
\end{proof}

\end{document}
